
qu'un
nombre fini de points singuliers, r\'eduction justifi\'ee par la formule de projection X 4.4.4 et par 4.5 (ii). Il existe donc une courbe ferm\'ee int\`egre $Y$ sur $X$ telle que $\langle c_1(L)\cdot Y\rangle\ne0$, et $Y$ est localement principale en dehors d'un ensemble fini $Z$. Soit $f:X'\to X$ l'\'eclatement de l'id\'eal $I$ de ${O}_X$. L'id\'eal $M=I\cdot{O}_{X'}$ est inversible, et il d\'efinit un sous-sch\'ema de $X'$, qui est la r\'eunion de deux sous-sch\'emas $Y'$, $Y''$ tels que $f(Y')=Y$ et $f(Y'')\subset Z$. Gr\^ace \`a la formule de projection X 4.4.4, on a que $\langle c_1(f^*L)\cdot c_1(M)\rangle=-\langle c_1(f^*L)\cdot Y'\rangle-\langle c_1(f^*L)\cdot Y''\rangle$ est \'egal \`a $-\langle c_1(L)\cdot Y\rangle-0\ne0$, et on peut remplacer $X,L,H$ par $X',f^*L,f^*H$ ; on est donc ramen\'e au cas pr\'ec\'edent, et la d\'emonstration de 7.1 est achev\'ee.

\textbf{Remarque 7.2.} On ignore si une surface irr\'eductible, propre (sur un corps alg\'ebriquement clos), qui poss\`ede un faisceau inversible $H$ tel que $\langle c_1(H)^2\rangle>0$, est n\'ecessairement projective. On ignore \'egalement si pour toute surface irr\'eductible, propre, la forme d'intersection sur $A^1_{\mathrm{num}}(X)_{\Q}$ est non-d\'eg\'en\'er\'ee ; il r\'esulte de 7.1 que lorsqu'elle est d\'eg\'en\'er\'ee, elle est n\'ecessairement n\'egative (non-d\'efinie).

\textbf{Corollaire 7.3.} \emph{Sous les conditions de 7.1, soit $M$ un faisceau inversible sur $X$ qui n'est pas num\'eriquement proportionnel \`a $H$, et soient $d=\langle c_1(M),c_1(H)\rangle$ et $h=\langle c_1(H)^2\rangle$. Alors, $\langle c_1(M)^2\rangle<d^2/h$.}

En effet, soit $L=M^{\otimes h}\otimes H^{\otimes -d}$. Comme $\langle c_1(L)\cdot c_1(H)\rangle=0$, d'apr\`es 7.1, $\langle c_1(L)^2\rangle<0$ ; d'o\`u l'assertion.

